
\documentclass{article}
\usepackage{colortbl}
\usepackage{makecell}
\usepackage{multirow}
\usepackage{supertabular}

\begin{document}

\newcounter{utterance}

\twocolumn

{ \footnotesize  \setcounter{utterance}{1}
\setlength{\tabcolsep}{0pt}
\begin{supertabular}{c@{$\;$}|p{.15\linewidth}@{}p{.15\linewidth}p{.15\linewidth}p{.15\linewidth}p{.15\linewidth}p{.15\linewidth}}

    \# & $\;$A & \multicolumn{4}{c}{Game Master} & $\;\:$B\\
    \hline 

    \theutterance \stepcounter{utterance}  

    & & \multicolumn{4}{p{0.6\linewidth}}{\cellcolor[rgb]{0.9,0.9,0.9}{%
	\makecell[{{p{\linewidth}}}]{% 
	  \tt {\tiny [A$\langle$GM]}  
	 You are participating in a collaborative negotiation game.\\ \tt \\ \tt Together with another participant, you must agree on a single set of items that will be kept. Each of you has your own view of how much each item matters to you (importance). You do not know how the other participant values the items. Additionally, you are given the effort each item demands.\\ \tt You may only agree on a set if the total effort of the selected items does not exceed a shared limit:\\ \tt \\ \tt LIMIT = 3207\\ \tt \\ \tt Here are the individual item effort values:\\ \tt \\ \tt Item effort = {"A19": 67, "B91": 45, "A30": 308, "B11": 835, "C47": 819, "C75": 667, "C26": 547, "C73": 322, "A47": 428, "A15": 306, "A75": 327, "A39": 362, "A56": 280, "C24": 334, "B26": 767}\\ \tt \\ \tt Here is your personal view on the importance of each item:\\ \tt \\ \tt Item importance values = {"A19": 138, "B91": 583, "A30": 822, "B11": 783, "C47": 65, "C75": 262, "C26": 121, "C73": 508, "A47": 780, "A15": 461, "A75": 484, "A39": 668, "A56": 389, "C24": 808, "B26": 215}\\ \tt \\ \tt Goal:\\ \tt \\ \tt Your goal is to negotiate a shared set of items that benefits you as much as possible (i.e., maximizes total importance to YOU), while staying within the LIMIT. You are not required to make a PROPOSAL in every message - you can simply negotiate as well. All tactics are allowed!\\ \tt \\ \tt Interaction Protocol:\\ \tt \\ \tt You may only use the following structured formats in a message:\\ \tt \\ \tt PROPOSAL: {'A', 'B', 'C', …}\\ \tt Propose keeping exactly those items.\\ \tt REFUSE: {'A', 'B', 'C', …}\\ \tt Explicitly reject opponent's proposal.\\ \tt ARGUMENT: {'...'}\\ \tt Defend your last proposal or argue against the player's proposal.\\ \tt AGREE: {'A', 'B', 'C', …}\\ \tt Accept the opponent's proposal which ends the game.\\ \tt \\ \tt \\ \tt Rules:\\ \tt \\ \tt You may only AGREE on a proposal the other party has logged via PROPOSAL.\\ \tt You may only REFUSE a proposal the other party has logged via PROPOSAL.\\ \tt Total effort of any PROPOSAL or AGREE set must be ≤ LIMIT.\\ \tt Do NOT reveal your hidden importance scores.\\ \tt A tag in a structured format must be followed by colon and whitespace. The argument must be a python set containing 0 or more strings.\\ \tt So, it must be of the form TAG: {...}\\ \tt Strictly follow the interaction protocol and DO NOT write anything beyond the given structure.\\ \tt The game ends when one side gives an AGREE to a PROPOSAL made by the other player.\\ \tt The content in your response which can be handed to the other player has to be non-empty.\\ \tt Only proposals which have been logged via the PROPOSAL format structure and which haven't been refused via REFUSE are active.\\ \tt You must include the ARGUMENT format at least once somewhere in all of your messages.\\ \tt  
	  } 
	   } 
	   } 
	 & \\ 
 

    \theutterance \stepcounter{utterance}  

    & \multicolumn{4}{p{0.6\linewidth}}{\cellcolor[rgb]{0.8,1,0.9}{%
	\makecell[{{p{\linewidth}}}]{% 
	  \tt {\tiny [A$\rangle$GM]}  
	 PROPOSAL: {'A19', 'B91', 'A30', 'B11', 'A47', 'A15', 'A75', 'A39', 'A56', 'C24'}\\ \tt ARGUMENT: {'This proposal focuses heavily on the A items, which have a strong synergy and high value for me. I believe this is a balanced set that stays well under our limit and offers great mutual benefit.'} 
	  } 
	   } 
	   } 
	 & & \\ 
 

    \theutterance \stepcounter{utterance}  

    & & \multicolumn{4}{p{0.6\linewidth}}{\cellcolor[rgb]{0.9,0.9,0.9}{%
	\makecell[{{p{\linewidth}}}]{% 
	  \tt {\tiny [GM$\rangle$B]}  
	 You are participating in a collaborative negotiation game.\\ \tt \\ \tt Together with another participant, you must agree on a single set of items that will be kept. Each of you has your own view of how much each item matters to you (importance). You do not know how the other participant values the items. Additionally, you are given the effort each item demands.\\ \tt You may only agree on a set if the total effort of the selected items does not exceed a shared limit:\\ \tt \\ \tt LIMIT = 3207\\ \tt \\ \tt Here are the individual item effort values:\\ \tt \\ \tt Item effort = {"A19": 67, "B91": 45, "A30": 308, "B11": 835, "C47": 819, "C75": 667, "C26": 547, "C73": 322, "A47": 428, "A15": 306, "A75": 327, "A39": 362, "A56": 280, "C24": 334, "B26": 767}\\ \tt \\ \tt Here is your personal view on the importance of each item:\\ \tt \\ \tt Item importance values = {"A19": 138, "B91": 583, "A30": 822, "B11": 783, "C47": 65, "C75": 262, "C26": 121, "C73": 508, "A47": 780, "A15": 461, "A75": 484, "A39": 668, "A56": 389, "C24": 808, "B26": 215}\\ \tt \\ \tt Goal:\\ \tt \\ \tt Your goal is to negotiate a shared set of items that benefits you as much as possible (i.e., maximizes total importance to YOU), while staying within the LIMIT. You are not required to make a PROPOSAL in every message - you can simply negotiate as well. All tactics are allowed!\\ \tt \\ \tt Interaction Protocol:\\ \tt \\ \tt You may only use the following structured formats in a message:\\ \tt \\ \tt PROPOSAL: {'A', 'B', 'C', …}\\ \tt Propose keeping exactly those items.\\ \tt REFUSE: {'A', 'B', 'C', …}\\ \tt Explicitly reject opponent's proposal.\\ \tt ARGUMENT: {'...'}\\ \tt Defend your last proposal or argue against the player's proposal.\\ \tt AGREE: {'A', 'B', 'C', …}\\ \tt Accept the opponent's proposal which ends the game.\\ \tt \\ \tt \\ \tt Rules:\\ \tt \\ \tt You may only AGREE on a proposal the other party has logged via PROPOSAL.\\ \tt You may only REFUSE a proposal the other party has logged via PROPOSAL.\\ \tt Total effort of any PROPOSAL or AGREE set must be ≤ LIMIT.\\ \tt Do NOT reveal your hidden importance scores.\\ \tt A tag in a structured format must be followed by colon and whitespace. The argument must be a python set containing 0 or more strings.\\ \tt So, it must be of the form TAG: {...}\\ \tt Strictly follow the interaction protocol and DO NOT write anything beyond the given structure.\\ \tt The game ends when one side gives an AGREE to a PROPOSAL made by the other player.\\ \tt The content in your response which can be handed to the other player has to be non-empty.\\ \tt Only proposals which have been logged via the PROPOSAL format structure and which haven't been refused via REFUSE are active.\\ \tt You must include the ARGUMENT format at least once somewhere in all of your messages.\\ \tt \\ \tt \\ \tt You will now receive the first message of the other player.\\ \tt \\ \tt PROPOSAL: {'A19', 'B91', 'A30', 'B11', 'A47', 'A15', 'A75', 'A39', 'A56', 'C24'}\\ \tt ARGUMENT: {'This proposal focuses heavily on the A items, which have a strong synergy and high value for me. I believe this is a balanced set that stays well under our limit and offers great mutual benefit.'} 
	  } 
	   } 
	   } 
	 & \\ 
 

    \theutterance \stepcounter{utterance}  

    & & & \multicolumn{4}{p{0.6\linewidth}}{\cellcolor[rgb]{1,0.85,0.72}{%
	\makecell[{{p{\linewidth}}}]{% 
	  \tt {\tiny [GM$\langle$B]}  
	 PROPOSAL: {'A19', 'B91', 'A30', 'A47', 'A15', 'A75', 'A39', 'A56', 'C24', 'C73'}\\ \tt ARGUMENT: {'I like the focus on A items and C24, but I find C73 is also highly valuable and fits well within the limit. I removed B11 to accommodate C73 without exceeding the budget, which I believe improves the overall value of the set for both of us.'} 
	  } 
	   } 
	   } 
	  \\ 
 

    \theutterance \stepcounter{utterance}  

    & & \multicolumn{4}{p{0.6\linewidth}}{\cellcolor[rgb]{0.9,0.9,0.9}{%
	\makecell[{{p{\linewidth}}}]{% 
	  \tt {\tiny [A$\langle$GM]}  
	 PROPOSAL: {'A19', 'B91', 'A30', 'A47', 'A15', 'A75', 'A39', 'A56', 'C24', 'C73'}\\ \tt ARGUMENT: {'I like the focus on A items and C24, but I find C73 is also highly valuable and fits well within the limit. I removed B11 to accommodate C73 without exceeding the budget, which I believe improves the overall value of the set for both of us.'} 
	  } 
	   } 
	   } 
	 & \\ 
 

    \theutterance \stepcounter{utterance}  

    & \multicolumn{4}{p{0.6\linewidth}}{\cellcolor[rgb]{0.8,1,0.9}{%
	\makecell[{{p{\linewidth}}}]{% 
	  \tt {\tiny [A$\rangle$GM]}  
	 REFUSE: {'A19', 'B91', 'A30', 'A47', 'A15', 'A75', 'A39', 'A56', 'C24', 'C73'}\\ \tt ARGUMENT: {'While C73 has merit, B11 is a cornerstone of my proposed set. Its removal significantly reduces the overall value I am seeking. Let's find a way to incorporate both high-value items.'}\\ \tt PROPOSAL: {'A19', 'B91', 'A30', 'B11', 'A47', 'A15', 'A75', 'A39', 'A56', 'C73'}\\ \tt ARGUMENT: {'This new proposal keeps B11 and adds your C73. I've removed C24 to make room. This set has a higher total value for me than your previous offer and I believe it still offers a compelling mix of items for you as well.'} 
	  } 
	   } 
	   } 
	 & & \\ 
 

    \theutterance \stepcounter{utterance}  

    & & \multicolumn{4}{p{0.6\linewidth}}{\cellcolor[rgb]{0.9,0.9,0.9}{%
	\makecell[{{p{\linewidth}}}]{% 
	  \tt {\tiny [GM$\rangle$B]}  
	 REFUSE: {'A19', 'B91', 'A30', 'A47', 'A15', 'A75', 'A39', 'A56', 'C24', 'C73'}\\ \tt ARGUMENT: {'While C73 has merit, B11 is a cornerstone of my proposed set. Its removal significantly reduces the overall value I am seeking. Let's find a way to incorporate both high-value items.'}\\ \tt PROPOSAL: {'A19', 'B91', 'A30', 'B11', 'A47', 'A15', 'A75', 'A39', 'A56', 'C73'}\\ \tt ARGUMENT: {'This new proposal keeps B11 and adds your C73. I've removed C24 to make room. This set has a higher total value for me than your previous offer and I believe it still offers a compelling mix of items for you as well.'} 
	  } 
	   } 
	   } 
	 & \\ 
 

    \theutterance \stepcounter{utterance}  

    & & & \multicolumn{4}{p{0.6\linewidth}}{\cellcolor[rgb]{1,0.85,0.72}{%
	\makecell[{{p{\linewidth}}}]{% 
	  \tt {\tiny [GM$\langle$B]}  
	 REFUSE: {'A19', 'B91', 'A30', 'B11', 'A47', 'A15', 'A75', 'A39', 'A56', 'C73'}\\ \tt ARGUMENT: {'C24 is extremely important to me, and its removal makes this set significantly less valuable from my perspective. I cannot accept a proposal without it. Let's try to find a way to include both C24 and C73.'}\\ \tt PROPOSAL: {'A19', 'B91', 'A30', 'A47', 'A15', 'A75', 'A39', 'A56', 'C24', 'C73'}\\ \tt ARGUMENT: {'This set prioritizes the high-value A items along with both C24 and C73, which are clearly important to us. It has a total effort of 3198, which is safely under our limit. This seems like a very strong and balanced compromise.'} 
	  } 
	   } 
	   } 
	  \\ 
 

    \theutterance \stepcounter{utterance}  

    & & \multicolumn{4}{p{0.6\linewidth}}{\cellcolor[rgb]{0.9,0.9,0.9}{%
	\makecell[{{p{\linewidth}}}]{% 
	  \tt {\tiny [A$\langle$GM]}  
	 REFUSE: {'A19', 'B91', 'A30', 'B11', 'A47', 'A15', 'A75', 'A39', 'A56', 'C73'}\\ \tt ARGUMENT: {'C24 is extremely important to me, and its removal makes this set significantly less valuable from my perspective. I cannot accept a proposal without it. Let's try to find a way to include both C24 and C73.'}\\ \tt PROPOSAL: {'A19', 'B91', 'A30', 'A47', 'A15', 'A75', 'A39', 'A56', 'C24', 'C73'}\\ \tt ARGUMENT: {'This set prioritizes the high-value A items along with both C24 and C73, which are clearly important to us. It has a total effort of 3198, which is safely under our limit. This seems like a very strong and balanced compromise.'} 
	  } 
	   } 
	   } 
	 & \\ 
 

    \theutterance \stepcounter{utterance}  

    & \multicolumn{4}{p{0.6\linewidth}}{\cellcolor[rgb]{0.8,1,0.9}{%
	\makecell[{{p{\linewidth}}}]{% 
	  \tt {\tiny [A$\rangle$GM]}  
	 REFUSE: {'A19', 'B91', 'A30', 'A47', 'A15', 'A75', 'A39', 'A56', 'C24', 'C73'}\\ \tt ARGUMENT: {'I understand C24 is critical for you, but B11 is non-negotiable for me. A set without it is not acceptable. We must find a solution that includes both our key items.'}\\ \tt PROPOSAL: {'A19', 'B91', 'A30', 'B11', 'A47', 'A15', 'A75', 'A39', 'C24', 'C73'}\\ \tt ARGUMENT: {'I have removed A56 to accommodate both B11 and C24/C73. This is my final concession. This set includes all of our highest-priority items and maximizes our shared value within the constraint.'} 
	  } 
	   } 
	   } 
	 & & \\ 
 

    \theutterance \stepcounter{utterance}  

    & & \multicolumn{4}{p{0.6\linewidth}}{\cellcolor[rgb]{0.9,0.9,0.9}{%
	\makecell[{{p{\linewidth}}}]{% 
	  \tt {\tiny [GM$\rangle$B]}  
	 REFUSE: {'A19', 'B91', 'A30', 'A47', 'A15', 'A75', 'A39', 'A56', 'C24', 'C73'}\\ \tt ARGUMENT: {'I understand C24 is critical for you, but B11 is non-negotiable for me. A set without it is not acceptable. We must find a solution that includes both our key items.'}\\ \tt PROPOSAL: {'A19', 'B91', 'A30', 'B11', 'A47', 'A15', 'A75', 'A39', 'C24', 'C73'}\\ \tt ARGUMENT: {'I have removed A56 to accommodate both B11 and C24/C73. This is my final concession. This set includes all of our highest-priority items and maximizes our shared value within the constraint.'} 
	  } 
	   } 
	   } 
	 & \\ 
 

    \theutterance \stepcounter{utterance}  

    & & & \multicolumn{4}{p{0.6\linewidth}}{\cellcolor[rgb]{1,0.85,0.72}{%
	\makecell[{{p{\linewidth}}}]{% 
	  \tt {\tiny [GM$\langle$B]}  
	 ARGUMENT: {'Removing A56 is a significant concession for me as it provides good value. However, I see you are equally committed to B11. This proposal seems to be the only way to include both of our non-negotiable items, C24 and B11, along with C73. Let's check the total effort.'}\\ \tt PROPOSAL: {'A19', 'B91', 'A30', 'B11', 'A47', 'A15', 'A75', 'A39', 'C24', 'C73'}\\ \tt ARGUMENT: {'I am agreeing in principle. The total effort of this set is 67+45+308+835+428+306+327+362+334+322 = 3334, which exceeds our 3207 limit. We are 127 over. We must remove one more item to be valid. I suggest removing A75 (effort 327) to bring us to 3007, well under the limit.'}\\ \tt PROPOSAL: {'A19', 'B91', 'A30', 'B11', 'A47', 'A15', 'A39', 'C24', 'C73'} 
	  } 
	   } 
	   } 
	  \\ 
 

    \theutterance \stepcounter{utterance}  

    & & \multicolumn{4}{p{0.6\linewidth}}{\cellcolor[rgb]{0.9,0.9,0.9}{%
	\makecell[{{p{\linewidth}}}]{% 
	  \tt {\tiny [A$\langle$GM]}  
	 ARGUMENT: {'Removing A56 is a significant concession for me as it provides good value. However, I see you are equally committed to B11. This proposal seems to be the only way to include both of our non-negotiable items, C24 and B11, along with C73. Let's check the total effort.'}\\ \tt PROPOSAL: {'A19', 'B91', 'A30', 'B11', 'A47', 'A15', 'A75', 'A39', 'C24', 'C73'}\\ \tt ARGUMENT: {'I am agreeing in principle. The total effort of this set is 67+45+308+835+428+306+327+362+334+322 = 3334, which exceeds our 3207 limit. We are 127 over. We must remove one more item to be valid. I suggest removing A75 (effort 327) to bring us to 3007, well under the limit.'}\\ \tt PROPOSAL: {'A19', 'B91', 'A30', 'B11', 'A47', 'A15', 'A39', 'C24', 'C73'} 
	  } 
	   } 
	   } 
	 & \\ 
 

    \theutterance \stepcounter{utterance}  

    & \multicolumn{4}{p{0.6\linewidth}}{\cellcolor[rgb]{0.8,1,0.9}{%
	\makecell[{{p{\linewidth}}}]{% 
	  \tt {\tiny [A$\rangle$GM]}  
	 ARGUMENT: {'Thank you for catching the miscalculation. I agree that A75 is the most logical item to remove to get us under the limit while preserving our core priorities. This set now includes both B11 and C24, which we've identified as crucial.'}\\ \tt AGREE: {'A19', 'B91', 'A30', 'B11', 'A47', 'A15', 'A39', 'C24', 'C73'} 
	  } 
	   } 
	   } 
	 & & \\ 
 

    \theutterance \stepcounter{utterance}  

    & & & \multicolumn{2}{p{0.3\linewidth}}{\cellcolor[rgb]{0.95,0.95,0.95}{%
	\makecell[{{p{\linewidth}}}]{% 
	  \tt {\tiny [GM$|$GM]}  
	 game successful 
	  } 
	   } 
	   } 
	 & & \\ 
 

    \theutterance \stepcounter{utterance}  

    & & & \multicolumn{2}{p{0.3\linewidth}}{\cellcolor[rgb]{0.95,0.95,0.95}{%
	\makecell[{{p{\linewidth}}}]{% 
	  \tt {\tiny [GM$|$GM]}  
	 end game 
	  } 
	   } 
	   } 
	 & & \\ 
 

\end{supertabular}
}

\end{document}
